\noindent Третья часть собирает из нейронов схемы. Сначала --- общие типы сетей: прямые и рекуррентные, и роль торможения. Потом --- конкретные задачи, которые мозг решает, и алгоритмы, которые он для этого использует: фильтрация изображения, спектральный анализ звука, управление движением, ассоциативная память, обучение с подкреплением, тактирование.
